% =====================================================================
%  Section 7 -- Limitations
% =====================================================================

\section{Limitations}
\label{sec:limitations}

We acknowledge the following limitations.

\paragraph{Training data and cross-sensor transfer.} The main benchmark trains on FSSD alone (663 RGB images, none of them water-free). We evaluate cross-domain transfer on 250 RGB chips from Sen1Floods11 (Section~\ref{sec:results-ood}): every FSSD-trained model collapses to IoU $\approx 0.16$ on this Sentinel-2 optical benchmark, so the recommended configurations do \emph{not} transfer zero-shot across sensors and viewpoints. Within the close-range RGB regime, FSSD-only training costs a $\sim$0.42-point flood-IoU drop on high-altitude UAV imagery and 72--100\% false-alarm rates on water-free scenes (Section~\ref{sec:results-zeroshot}). Retraining on a combined corpus with water-free images removes most of that cost (false-alarm rates 3--9\%, Section~\ref{sec:results-combined}), but that experiment covers a reduced grid: no-KD and combined KD only, no MobileNetV2 baseline, and no latency re-measurement, so the per-platform latencies in Table~\ref{tab:t5-latency} were measured on FSSD-trained exports (the architectures, and hence the latencies, are identical). Its single model-selection fold is also small (54--55 images) and contains no water-free image, which may explain the teacher's unstable false-alarm rate. Cross-sensor deployment (e.g.\ to satellite optical or SAR) would require domain adaptation or in-sensor retraining, which we leave to future work, and two close-range training sources are not enough to claim broad geographic robustness.

\paragraph{RGB-only modality.} Operational disaster-response platforms increasingly fuse SAR (cloud-penetrating, all-weather) with optical imagery; our work is restricted to the RGB regime FSSD was constructed for. The two persistent failure modes in Section~\ref{sec:discussion-failures} (wet-asphalt confusion, vegetation occlusion) are intrinsic to the RGB modality and cannot be solved by a better RGB model.

\paragraph{Three folds at the fold level.} The 3-fold protocol gives a minimum reportable two-sided Wilcoxon p of 0.25 \emph{at the fold level}, so we do not rely on fold-level hypothesis tests. Significance for the KD contrasts is instead established by the per-image paired bootstrap over the pooled held-out validation set ($n=399$ images; Sections~\ref{sec:setup-stats},~\ref{sec:results-kd}), which has adequate power on the FSSD test distribution. Two caveats remain. First, the bootstrap quantifies uncertainty over \emph{images}, not over independent dataset draws, and the three fold training sets overlap by construction, so it does not license population-level generalisation claims --- a larger fold count or the independent-dataset evaluation flagged above would be required for that. Second, the PTQ collapses are reported as effect-size magnitudes (\eg the $-0.45$ IoU drop on MobileNetV3-S INT8 is unambiguous) rather than by hypothesis test, given their trimodal cross-fold behaviour (Table~\ref{tab:t4-quant}).

\paragraph{Two deployment platforms (M2 Pro + browser WASM).} We do not measure Raspberry~Pi~4, AWS Graviton (ARM CPU), Android (Termux $+$ TFLite), or any GPU edge target (Jetson Nano / Orin). Adding these is largely an artifact-conversion exercise --- the ONNX hub artifact remains the same --- but each adds a non-trivial setup and tuning cost we did not budget for. Pi~4 and Graviton numbers are the most directly comparable extensions.

\paragraph{Post-training quantization only.} We do not evaluate quantization-aware training (QAT). The 0.45-IoU PTQ collapse on MobileNetV3-Small and MobileViT-XXS is a candidate target for QAT --- both architectures are reportedly QAT-recoverable in the classification literature --- but QAT is an order-of-magnitude larger training investment than the experiments reported here and is left to future work.

\paragraph{No energy / power measurements.} We do not report Joules-per-inference. Energy is the metric most directly relevant to battery-operated deployments (drones, field tablets) and we acknowledge its absence; the equipment required (USB-meter, INA219 / current-clamp logger) was not available during this study.

\paragraph{Saturated benchmark.} FSSD saturates at IoU $\approx 0.93$ for the entire UNet $+$ ImageNet-pretrained backbone family. The negative-KD finding is therefore strictly a claim about saturated-benchmark behaviour. On the less saturated combined corpus (Section~\ref{sec:results-combined}) combined KD does recover a flood-IoU gain, at the cost of inheriting the teacher's false alarms. That experiment uses a single teacher per fold and does not separate the response and feature terms, so which term carries the inherited errors remains open.
